\subsection{The SIR model}
The classical SIR model partitions a closed population of size $N$ into susceptible ($S$),
infectious ($I$), and recovered ($R$) compartments:
\begin{align}
  \frac{dS}{dt} &= -\frac{\beta S I}{N}, &
  \frac{dI}{dt} &= \frac{\beta S I}{N} - \gamma I, &
  \frac{dR}{dt} &= \gamma I. \label{eq:sir}
\end{align}
The basic reproduction number is $R_0 = \beta/\gamma$: an epidemic grows while $R_0 \cdot S(t)/N > 1$
and declines once susceptible depletion pushes it below 1.

\subsection{Base paper and our extension}
Capaldi et al.~\cite{capaldi2012} fit \eqref{eq:sir} to synthetic and real incidence data via
least squares, quantify the uncertainty in the resulting $(\hat\beta,\hat\gamma)$ via an asymptotic
covariance estimate, and study how estimation quality degrades with sparser sampling. Their solver
is used only as a means to an end. We take their estimation machinery as a starting point but
restructure the experiment so that \emph{solver identity and step size} are first-class independent
variables, crossed with noise level, sampling interval, observation window, and the underlying
$(\beta,\gamma)$ regime.

\subsection{Related numerical-methods reference}
Prodanov~\cite{prodanov2021} gives an analytical treatment of the SIR system's phase-plane
reduction that we build on directly for our gold-standard solution (Section~\ref{sec:methods}).
